\documentclass[12pt,a4paper]{article}
\usepackage[utf8]{inputenc}
\usepackage[T1]{fontenc}
\usepackage{amsmath,amssymb,amsthm}
\usepackage{geometry}
\usepackage{longtable}
\usepackage{booktabs}
\usepackage{array}
\usepackage{multicol}
\usepackage{fancyhdr}
\usepackage{titlesec}
\usepackage{enumitem}
\usepackage{mathrsfs}

\geometry{margin=1in}

\pagestyle{fancy}
\fancyhf{}
\rhead{Cayley --- Collected Papers, Vol.~XIII}
\lhead{Pages 451--500}
\cfoot{\thepage}

\titleformat{\section}{\large\bfseries}{\thesection.}{0.5em}{}
\titleformat{\subsection}{\normalsize\bfseries}{\thesubsection.}{0.5em}{}

\newcommand{\primd}{\textit{prim\`a facie}}

\begin{document}

\begin{center}
{\Large\bfseries The Collected Mathematical Papers of\\[0.5em] Arthur Cayley, Vol.~XIII}\\[1em]
{\large Pages 451--500}\\[0.5em]
{\normalsize Transcription}
\end{center}

\hrule
\vspace{1em}


\section*{947.}
\addcontentsline{toc}{section}{947. On a System of Two Tetrads of Circles; and Other Systems of Two Tetrads}

\begin{center}
{\large\bfseries ON A SYSTEM OF TWO TETRADS OF CIRCLES; AND OTHER SYSTEMS OF TWO TETRADS.}
\end{center}

\begin{flushright}
[From the Proceedings of the Cambridge Philosophical Society, vol.~VIII. (1893),\\
pp.~54---59.]
\end{flushright}

The investigations of the present paper were suggested to me by Mr Orr's paper, ``The Contacts of certain Systems of Circles,'' Proc.~Camb.~Phil.~Soc., vol.~vii.

\textbf{1.} It is possible to find in piano two tetrads of circles, or say four red circles and four blue circles, such that each red circle touches each blue circle: in fact, counting the constants, a circle depends upon 3 constants, or say it has a capacity $=3$; the capacity of the eight circles is thus $=24$; and the postulation or number of conditions to be satisfied is $=16$: the resulting capacity of the system is thus \primd, $16-24=8$. It will, however, appear that, in the system considered, the true value is $=9$.

\textbf{2.} The \primd\ value of the capacity being $=8$, we are not at liberty to assume at pleasure three circles of the system. And, in fact, assuming at pleasure say 3 red circles, then touching each of these we have 8 circles, forming $\frac{1}{24}\, 8 \cdot 7 \cdot 6 \cdot 5 = 70$, tetrads of circles: taking at random any one of these tetrads for the blue circles, the remaining red circle has to be determined so as to touch each of the four blue circles, that is, by four instead of three conditions; and there is not in general any red circle satisfying these four conditions. But the 8 tangent circles do not stand to each other in a relation of symmetry, but form in fact four pairs of circles; and it is possible out of the 70 tetrads to select (and that in 6 ways) a tetrad of blue circles, such that there exists a fourth red circle touching each of these four blue circles. We have thus a system depending upon 3 arbitrary circles, and for which, therefore, the capacity is $=9$. It is (as is known) possible, in quite a different manner, out of the 70 tetrads to select (and that in 8 ways) a tetrad of blue circles such that there exists a fourth red circle touching each of these four blue circles---but the present paper relates exclusively to the first-mentioned 6 tetrads and not to these 8 tetrads.

\textbf{3.} I consider, in the first instance, a particular case in which the three red circles are not all of them arbitrary, but have a capacity $9-1=8$; and pass from this to the general case where the capacity is $=9$. Calling the red circles 1, 2, 3 and 4; I start with the circles 1 and 2 arbitrary, and 3 a circle equal to 2: the radical axis, or common chord, of the circles 2 and 3 is thus a line bisecting at right angles the line joining the centres of the circles 2 and 3, say this is the line~$\Omega$. We have then four circles, each having its centre in the line~$\Omega$ and touching the circles 1 and 2: in fact, the locus of the centre of a circle touching the circles 1 and 2 is a pair of conics, each of them having for foci the centres of these circles: the line~$\Omega$ meets each of these conics in two points, and there are thus on the line~$\Omega$ four points, each of them the centre of a circle touching the circles 1 and 2. But the equal circles 2 and 3 are symmetrically situate in regard to the line~$\Omega$; and it is obvious that the four circles, having their centres on the line~$\Omega$, will each of them also touch the circle~3; we have thus the four blue circles, each of them with its centre on the line~$\Omega$ and touching each of the red circles 1, 2 and 3. And it is moreover clear that, taking the red circle 4 equal to 1 and situate symmetrically therewith in regard to the line~$\Omega$, then this circle 4 will touch each of the blue circles: so that we have here the four blue circles, each of them touching the four red circles. As already mentioned, the blue circles have their centre on the line~$\Omega$, that is, the line~$\Omega$ is a common orthotomic of the four blue circles.

\textbf{4.} By inverting in regard to an arbitrary circle we pass to the general case; the line~$\Omega$ becomes thus a circle~$\Omega$, orthotomic to each of the blue circles. Starting \textit{ab initio}, we have here at pleasure the red circles 1, 2, 3: the circle $\Omega$ is a circle having for centre a centre of symmetry of the circles 2 and 3, and passing through the points of intersection (real or imaginary) of these two circles; the circles 2 and 3 are thus the inverses (or say the images) each of the other in regard to the circle~$\Omega$. We can then find 4 circles each of them orthotomic to $\Omega$, and touching the circles 1 and 2: but a circle orthotomic to $\Omega$ is its own inverse or image in regard to $\Omega$; and it will thus touch the circle 3 which is the image of 2 in regard to $\Omega$. We have thus the four blue circles each of them touching the red circles 1, 2 and 3; and then, taking the red circle 4 as the inverse or image of 1 in regard to $\Omega$, this circle 4 will also touch each of the blue circles. Thus starting with the arbitrary red circles 1, 2, 3, we find the four blue circles and the remaining red circle 4, such that each of the blue circles touches each of the red circles. Since in the construction we group together at pleasure the two circles 2, 3 (out of the three circles 1, 2, 3) and use at pleasure either of the two centres of symmetry, it appears that the number of ways in which the figure might have been completed is $=6$.

\textbf{5.} The blue circles have a common orthotomic circle $\Omega$, that is, the radical axis or common chord of each two of the blue circles passes through one and the same point, the centre of the circle $\Omega$. The figure is symmetrical in regard to the red and blue circles respectively, and thus the red circles have a common orthotomic circle $\Omega'$, that is, the radical axis or common chord of each two of the red circles passes through one and the same point, the centre of the circle~$\Omega'$.

\textbf{6.} Projecting stereographically on a spherical surface, the four red circles and the four blue circles become circles of the sphere; and then making the general homographic transformation, they become plane sections of a quadric surface; we have thus the theorem that on a given quadric surface it is possible to find four red sections and four blue sections such that each blue section touches each red section; and moreover the capacity of the system is $=9$; viz.~3 of the red sections may be assumed at pleasure. But (as is well known) the theory of the tangency of plane sections of a quadric surface is far more simple than that of the tangency of circles: the condition in order that two sections may touch each other is simply the condition that the line of intersection of the two planes shall touch the quadric surface. And we construct, as follows, the sections touching each of three given sections: say the given sections are 1, 2, 3; through the sections 1 and 2 we have two quadric cones having for vertices say the points $d_{12}$ and $i_{12}$ (direct and inverse centres of the two sections): similarly, through the sections 1 and 3 we have two quadric cones vertices $d_{13}$ and $i_{13}$ respectively, and through the sections 2 and 3 we have two quadric cones vertices $d_{23}$ and $i_{23}$ respectively; the points $d_{12}$, $i_{12}$, $d_{13}$, $i_{13}$, $d_{23}$, $i_{23}$ lie three and three in four intersecting lines or axes, viz. these are $d_{12}d_{13}d_{23}$, $d_{12}i_{13}i_{23}$, $d_{13}i_{12}i_{23}$, $d_{23}i_{12}i_{13}$ respectively. Through any one of these axes, say $d_{12}d_{13}d_{23}$, we may draw to the quadric surface two tangent planes each touching the three cones which have their vertices in the points $d_{23}$, $d_{31}$, $d_{12}$ respectively; and the section by either tangent plane is thus a section touching each of the three given sections 1, 2, 3; we have thus the eight tangent sections of these three sections.

\textbf{7.} Taking as three of the red sections the arbitrary sections 1, 2, 3; and grouping together two at pleasure of these sections, say 2 and 3; we may take for the blue sections the two sections through the axis $d_{23}d_{31}d_{12}$, and those through the axis $d_{23}i_{31}i_{12}$; we have thus the four blue sections touching each of the given red sections 1, 2, 3; and this being so, there exists a remaining red section 4 touching each of the blue sections; we have thus the four blue sections touching each of the red sections 1, 2, 3 and 4. This implies that the vertices or points $d_{24}$ and $d_{34}$ lie on the axis $d_{12}d_{13}d_{23}$, and that the vertices or points $i_{24}$ and $i_{34}$ lie on the axis $d_{23}i_{24}i_{13}$; or, what is the same thing, that the four sections 1, 2, 3, 4 have in common an axis $d_{12}d_{13}d_{14}d_{23}$ and also an axis $d_{23}i_{24}i_{13}i_{34}$.

\textbf{8.} If the quadric surface be a flat surface (surface aplatie) or conic, then the red sections become chords of the conic; the axes are lines in the plane of the conic, and thus the tangent planes through an axis each coincide with the plane of the conic, and it would at first sight appear that any theorem as to tangency becomes nugatory. But this is not so; comparing with the last preceding paragraph, we still have the theorem: on a given conic, taking at pleasure any three chords 1, 2, 3, it is possible to find a fourth chord 4, such that the four chords have in common an axis $d_{12}d_{13}d_{14}d_{23}$ and also an axis $d_{23}i_{24}i_{13}i_{34}$. And the analytical theory (although somewhat complex) is extremely interesting. Considering the conic $xz-y^{2}=0$, the coordinates of a point on the conic are given by $x:y:z = 1:\theta:\theta^{2}$, or say any point of the conic is determined by its parameter $\theta$; and this being so, considering any three chords 1, 2, 3, I take for the two extremities of 1 the values $\epsilon$, $\zeta$; for those of 2 the values $\alpha$, $\beta$; and for those of 3 the values $\gamma$, $\delta$; the remaining chord 4 is to be determined as above, and I take for its two extremities the values $E$, $Z$.

\textbf{9.} Starting with the chords 1, 2, 3, we have each of the points $d_{23}$, \&c., as the intersection of two lines, viz. these are
\[
\begin{aligned}
&\alpha z - y(\alpha+\delta) + z = 0, \quad \alpha z - y(\alpha+\gamma) + z = 0, \quad \text{for } d_{23}, i_{23},\\
&\alpha z - y(\alpha+\zeta) + z = 0, \quad \alpha z - y(\alpha+\epsilon) + z = 0, \quad \text{for } d_{12}, i_{12},\\
&\gamma z - y(\gamma+\epsilon) + z = 0, \quad \gamma z - y(\gamma+\zeta) + z = 0, \quad \text{for } d_{31}, i_{31},
\end{aligned}
\]
and we thence find without difficulty for the axis $d_{23}d_{31}d_{12}$ the equation
\begin{multline*}
a\{\alpha\beta(\gamma\delta - \zeta\epsilon) + \gamma\delta(\zeta\alpha - \epsilon\beta) + \zeta\epsilon(\beta\gamma - \alpha\delta)\} \\
+ y\{(\beta-\alpha)(\gamma\delta - \zeta\epsilon) + (\delta-\gamma)(\zeta\alpha - \epsilon\beta) + (\zeta-\epsilon)(\alpha\beta - \gamma\delta)\} \\
+ z\{\quad\quad - (\gamma\delta - \zeta\epsilon) - (\zeta\alpha - \epsilon\beta) - (\alpha\beta - \gamma\delta)\} = 0;
\end{multline*}
the equation of the axis $d_{23}i_{12}i_{13}$ is obtained herefrom by the interchange of $\epsilon$ and $\zeta$.

\textbf{10.} The points $d_{24}$ and $i_{34}$ will lie upon the first-mentioned axis if only $d_{24}$ lies upon this axis, viz. if we have
\begin{multline*}
\tfrac{1}{2}\{(\beta-\alpha)(\gamma\delta - \zeta\epsilon) + (\delta-\gamma)(\zeta\alpha - \epsilon\beta) + (\zeta-\epsilon)(\alpha\beta - \gamma\delta)\}\,(\epsilon E - \zeta) \\
+ \{\quad\quad - (\gamma\delta - \zeta\epsilon) - (\zeta\alpha - \epsilon\beta) - (\beta\gamma - \alpha\delta)\}\,\dotfill
\end{multline*}
Reducing this equation, the factor $\alpha - \beta$ divides out, and we finally obtain
\[
(\beta-\delta)(\alpha\gamma - \epsilon E) - (\alpha\delta - \beta\gamma)(\zeta - E) + (\alpha-\gamma)(\beta\epsilon - \delta\zeta) = 0,
\]
say this is
\[ A + BE + CZ + DEZ = 0, \]
where
\[
\begin{aligned}
A &= \dotfill, \\
B &= -(\gamma-\alpha)\beta\delta \cdot -(\alpha\beta - \gamma\delta)\epsilon + (\beta\delta)\zeta, \\
C &= -(\beta-\delta)\alpha\gamma + (\alpha\beta - \gamma\delta)\epsilon \cdot \dotfill + (\gamma-\alpha)\dotfill, \\
D &= -(\alpha\delta - \beta\gamma) - (\beta-\delta)\epsilon - (\gamma-\alpha)\zeta;
\end{aligned}
\]
viz. this is the condition for the existence of the axis $d_{23}d_{13}d_{12}d_{24}d_{34}$.

We interchange herein $\epsilon$, $\zeta$ and also $E$, $Z$, and we thus obtain
\[ A' + C'E + B'Z + D'EZ = 0, \]
where
\[
\begin{aligned}
A' &= \dotfill \quad (\zeta\delta - \alpha)\beta\gamma + (\gamma-\alpha)\zeta\delta\epsilon + (\alpha\delta - \beta\gamma)\zeta, \\
B' &= (\alpha\zeta - \beta\gamma)\epsilon \cdot \dotfill + (\zeta-\delta)\dotfill, \\
C' &= \dotfill \cdot -(\alpha\beta - \gamma\delta)\zeta + (\gamma-\alpha)\dotfill, \\
D' &= \dotfill;
\end{aligned}
\]
viz. this is the condition for the existence of the axis $d_{23}i_{13}i_{12}i_{34}i_{24}$.

\textbf{11.} I remark that we have
\[
A' = A + a(\epsilon - \zeta), \quad B' = B + b(\epsilon - \zeta), \quad C' = C + c(\epsilon - \zeta), \quad D' = D + d(\epsilon - \zeta),
\]
where
\[
\begin{aligned}
a &= (\beta-\delta)\alpha\gamma - (\gamma-\alpha)\beta\delta = \alpha\beta(\gamma+\delta) - \gamma\delta(\alpha+\beta), \\
b &= c = \gamma\delta - \alpha\beta, \\
d &= \alpha + \beta - \gamma - \delta,
\end{aligned}
\]
and further that
\[
\dotfill
\]
say this is $\Pi$.

\textbf{12.} It thus appears that, for the determination of $E$, $Z$, we have
\[
\begin{aligned}
A + BE + CZ + DEZ &= 0, \\
A' + C'E + B'Z + D'EZ &= 0.
\end{aligned}
\]
Eliminating $Z$, we find
\[
\begin{vmatrix}
A + BE & C + DE \\
A' + C'E & B' + D'E
\end{vmatrix} = 0,
\]
that is,
\[
(AC' - A'C) + (AD' - A'D + BB' - CC')\,E + (BD' - B'D)\,E^{2} = 0;
\]
upon reducing the coefficients of this equation it appears that they contain each of them the factor $\Pi$, and throwing out this factor, the equation is
\[
\epsilon\{\alpha\beta(\gamma-\delta) + \gamma\delta(\beta-\alpha) + (\alpha\delta - \beta\gamma)\} + \dotfill\,E = 0,
\]
this contains obviously the factor $E - \epsilon$, or throwing out this factor, we have for $E$ the simple equation
\[ E = \dotfill \]
In a similar manner it may be shown that the two equations give for $Z$ the like simple equation
\[ Z = \dotfill \]
viz. starting from the chords 1, 2, 3 which depend on the parameters $(\epsilon, \zeta)$, $(\alpha, \beta)$, $(\gamma, \delta)$ respectively, these last two equations give the parameters $(E, Z)$ of the chord~4.

\vspace{1em}
\hrule
\vspace{1em}


\newpage
\section*{948.}
\addcontentsline{toc}{section}{948. Report on the Pellian Equation}

\begin{center}
{\large\bfseries REPORT OF A COMMITTEE APPOINTED FOR THE PURPOSE OF CARRYING ON THE TABLES CONNECTED WITH THE PELLIAN EQUATION FROM THE POINT WHERE THE WORK WAS LEFT BY DEGEN IN 1817.}
\end{center}

\begin{flushright}
[From the British Association Report, (1893), pp.~73---120.]
\end{flushright}

We have, on the Pellian Equation, Degen's tables, the title of which is ``Canon Pellianus sive Tabula simplicissimam aequationis celebratissimae $y^{2} = ax^{2} + 1$ solutionem pro singulis numeri dati valoribus ab 1 usque ad 1000 in numeris rationalibus iisdemque integris exhibens.'' Autore Carolo Ferdinando Degen. Hafniae, apud Gerardum Bonnierum, MDCCCXVII., 8vo. Introductio, pp.~v---xxiv. Tabula~I. Solutionem aequationis $y^{2} - ax^{2} - 1 = 0$ exhibens, pp.~3---106. Tabula~II. Solutionem aequationis $y^{2} - ax^{2} + 1 = 0$, quotiescunque valor ipsius $a$ talem admiserit, exhibens, pp.~109---112.

The mode of calculation is explained in the Introduction, and illustrated by the examples of the numbers 209, 173.

As to the first of these, the entry in Table~I. is
\begin{center}
\begin{tabular}{ccccc}
14, & 2, 5, 3, & (2) & & \\
1, & 13, 5, 8, & 11 & & \\
209 & & & & \\
& 3220 & & & \\
& & 46551 & &
\end{tabular}
\end{center}
where the first line gives the expression of $\surd 209$ as a continued fraction, viz. we have
\[
\surd 209 = 14 +
\cfrac{1}{2+
 \cfrac{1}{5+
  \cfrac{1}{3+
   \cfrac{1}{2+
    \cfrac{1}{3+
     \cfrac{1}{5+
      \cfrac{1}{2+
       \cfrac{1}{28+\cdots}}}}}}}}
\]
the denominators being 2, 5, 3, (2), 3, 5, 2, then 28, which is the double of the integer part 14, and then again 2, 5, 3, (2), 3, 5, 2, and so on, the parentheses of the (2) being used to indicate that this is the middle term of the period.

The second row gives auxiliary numbers occurring in the calculation of the first row and having a meaning, as will presently appear. Observe that the 11 which comes under the (2) should also be printed in parentheses (11), but this is not done.

The process for the calculation of the $x$, $y$ is as follows:

\begin{center}
\begin{tabular}{r|r|r|r}
209 & & & \\ \hline
14 & & & \\
1 & 0 & $+1$ & \\
2 & 14 & 1 & $-13$ \\
5 & 29 & 2 & $+5$ \\
3 & 159 & 11 & $-8$ \\
(2) & 506 & 35 & $+(11)$ \\
3 & 1171 & 81 & $-8$ \\
5 & 4019 & 278 & $+5$ \\
2 & 21266 & 1471 & $-13$ \\
28 & 46551 & 3220 & $+1$
\end{tabular}
\end{center}

viz. writing down as a first column the numbers of the first row, and beginning the second column with 1, 14 (14 the number at the head of the first column), and the third column with 0, 1, we calculate the numbers of the second column, $29 = 2 \cdot 14 + 1$, $159 = 5 \cdot 29 + 14$, $506 = 3 \cdot 159 + 29$, \&c., and the numbers of the third column in like manner, $2 = 2 \cdot 1 + 0$, $11 = 5 \cdot 2 + 1$, $35 = 3 \cdot 11 + 2$, \&c.; and then writing down as a fourth column the numbers of the second row with the signs $+$, $-$ alternately, we have a series of equations $y^{2} - ax^{2} = \pm A$, viz.
\[
\begin{aligned}
1^{2} - 209 \cdot 0^{2} &= +1, \\
14^{2} - 209 \cdot 1^{2} &= -13, \\
29^{2} - 209 \cdot 2^{2} &= +5,
\end{aligned}
\]
the last of them being
\[ (46551)^{2} - 209 \cdot (3220)^{2} = +1, \]
this last corresponding as above to the value $+1$, and the numbers 46551 and 3220 being accordingly the $y$ and $x$ given in the fourth and third rows of the table.


\newpage
As to the second of the foregoing numbers, 173, the only difference is that the period has a double middle term, viz. the entry in the Table~I. is
\begin{center}
\begin{tabular}{cccc}
13, & 6, & (\phantom{1}1, & \phantom{1}1) \\
1, & 4, & (13, & 13) \\
190060 & & & \\
2499849 & & &
\end{tabular}
\end{center}
The first row gives the expression of $\surd 173$, viz. that is
\[
\surd 173 = 13 +
\cfrac{1}{6+
 \cfrac{1}{1+
  \cfrac{1}{1+
   \cfrac{1}{6+
    \cfrac{1}{26+\cdots}}}}
}
\]
the denominators being 6, 1, 1, 6, then 26 (the double of the integer part 13), and then again 6, 1, 1, 6, and so on. In the second row I remark that Degen prints the parentheses (13,~13) for the double middle term.

The process for the calculation of the $x$, $y$ is similar to that in the former case, viz. we have

\begin{center}
\begin{tabular}{r|r|r|r}
173 & & & \\ \hline
13 & & & \\
1 & 0 & $+1$ & \\
6 & 13 & 1 & $-4$ \\
1 & 79 & 6 & $+13$ \\
1 & 92 & 7 & $-13$ \\
6 & 171 & 13 & $+4$ \\
26 & 1118 & 85 & $-1$
\end{tabular}
\end{center}

where the second and third columns begin 1, 13 and 0, 1 respectively, and the remaining terms are calculated $79 = 6 \cdot 13 + 1$, $92 = 1 \cdot 79 + 13$, \&c., and $6 = 6 \cdot 1 + 0$, $7 = 1 \cdot 6 + 1$, \&c.; and then writing down as a fourth column the terms of the second row with the signs $+$, $-$ alternately, we have
\[
\begin{aligned}
1^{2} - 173 \cdot 0^{2} &= +1, \\
13^{2} - 173 \cdot 1^{2} &= -4, \\
79^{2} - 173 \cdot 6^{2} &= +13,
\end{aligned}
\]
the last equation being
\[ (1118)^{2} - 173 \cdot (85)^{2} = -1. \]

\newpage
\noindent\textbf{[948}
\addcontentsline{toc}{subsection}{948 (continued). Specimen Table}

The term for the last equation being always in a case such as the present one, not $+1$, but $-1$. The final numbers 1118, 85 are consequently entered not in Table~I., but in Table~II., viz. the entry in this table is
\begin{center}
\begin{tabular}{c}
173 \\
85 \\
1118
\end{tabular}
\end{center}
and thence we calculate the numbers $y$, $x$ of Table~I., viz. these are
\[
\begin{aligned}
2499849 &= 2 \cdot (1118)^{2} + 1, \\
190060 &= 2 \cdot 1118 \cdot 85.
\end{aligned}
\]

Generally Table~II. gives for each value of $a$, comprised therein, values of $x$, $y$, such that $y^{2} = ax^{2} - 1$, and then writing $y_{1} = 2y^{2} + 1$, $x_{1} = 2xy$, we have
\[
y_{1}^{2} = (2ax^{2} - 1)^{2} = 4a^{2}x^{4} - 4ax^{2} + 1 = a \cdot 4x^{2}(ax^{2} - 1) + 1 = ax_{1}^{2} + 1,
\]
so that $x_{1}$, $y_{1}$ are for the same value of $a$ the values of $x$, $y$ in Table~I.

It is to be remarked that the heading of Table~II. is not perfectly accurate, for it purports to give for every value of $a$, for which a solution exists, a solution of the equation $y^{2} = ax^{2} - 1$. What it really gives is the solution for each value of $a$ for which the period has a double middle term. But if $a = \alpha^{2} + 1$, then obviously we have a solution $y = \alpha$, $x = 1$, and for any such value of $a$ the period has a single middle term, viz. the entry in Table~I. is

\begin{center}
\begin{tabular}{cc}
$\alpha^{2}+1$ & \\
$\alpha$, & $(2\alpha)$ \\
1, & 1 \\
$2\alpha$ & \\
$2\alpha^{2}+1$ &
\end{tabular}
\end{center}

and we, in fact, have

\begin{center}
\begin{tabular}{r|r|r|r}
$\alpha^{2}+1$ & & & \\ \hline
$\alpha$ & & & \\
1 & 0 & $+1$ & \\
$(2\alpha)$ & $\alpha$ & 1 & $-1$ \\
$2\alpha$ & $2\alpha^{2}+1$ & $2\alpha$ & $+1$
\end{tabular}
\end{center}
\[
\begin{aligned}
1^{2} - (\alpha^{2}+1) \cdot 0^{2} &= +1, \\
\alpha^{2} - (\alpha^{2}+1) \cdot 1^{2} &= -1, \\
(2\alpha^{2}+1)^{2} - (\alpha^{2}+1) \cdot (2\alpha)^{2} &= +1.
\end{aligned}
\]

\newpage
The foregoing instances of the calculation of $x$, $y$ in the case of the numbers 209 and 173 suggest a table which may be regarded as an extended form of Degen's tables; viz. such a table, from $a = 2$ to $a = 99$, is as follows:

\begin{center}
\textbf{SPECIMEN OF EXTENDED FORM OF TABLE IN REGARD TO THE PELLIAN EQUATION.}
\end{center}

\small
\begin{longtable}{|c|c|c|c||c|c|c|c|}
\hline
$a$ & $y$ & $x$ & $y^{2} - ax^{2}$ & $a$ & $y$ & $x$ & $y^{2} - ax^{2}$ \\
\hline
\endfirsthead
\hline
$a$ & $y$ & $x$ & $y^{2} - ax^{2}$ & $a$ & $y$ & $x$ & $y^{2} - ax^{2}$ \\
\hline
\endhead
\hline
\endfoot
2 & 1 & 0 & $+1$ & 13 & 3 & 1 & $+1$ \\
  & (2) & 1 & 1 & & (1) & 1 & $-4$ \\
  & 3 & 2 & $+1$ & & 4 & 1 & $+3$ \\
\hline
3 & 1 & 0 & $+1$ & 14 & 3 & 1 & $+1$ \\
  & (1) & 1 & $-2$ & & (1) & 1 & $-5$ \\
  & 2 & 1 & $+1$ & & 4 & 1 & $+2$ \\
\hline
5 & 2 & 1 & $+1$ & 15 & 4 & 1 & $+1$ \\
  & (4) & 1 & $-1$ & (1) & 1 & $-14$ \\
  & 9 & 4 & $+1$ & 1 & 4 & $+1$ \\
\hline
6 & 2 & 1 & $+1$ & 17 & 4 & 1 & $+1$ \\
  & (2) & 2 & $-2$ & & (8) & 1 & $-1$ \\
  & 5 & 2 & $+1$ & & 33 & 8 & $+1$ \\
\hline
7 & 2 & 1 & $+1$ & 18 & 4 & 1 & $+1$ \\
  & (1) & 1 & $-3$ & & (4) & 1 & $+2$ \\
  & 3 & 1 & $+2$ & & 17 & 4 & $+1$ \\
\hline
8 & 2 & 1 & $+1$ & 19 & 4 & 1 & $+1$ \\
  & (1) & 1 & $-4$ & & (2) & 3 & $-5$ \\
  & 3 & 1 & $+1$ & & 39 & 9 & $+1$ \\
\hline
10 & 3 & 1 & $+1$ & 20 & 4 & 1 & $+1$ \\
   & (6) & 1 & $-1$ & & (2) & 1 & $-4$ \\
   & 19 & 6 & $+1$ & & 9 & 2 & $+1$ \\
\hline
11 & 3 & 1 & $+1$ & 21 & 4 & 1 & $+1$ \\
   & (3) & 1 & $-2$ & & (1) & 2 & $+3$ \\
   & 10 & 3 & $+1$ & & 5 & 1 & $+4$ \\
\hline
12 & 3 & 1 & $+1$ & 22 & 4 & 1 & $+1$ \\
   & (2) & 1 & $-3$ & & (1) & 2 & $-6$ \\
   & 7 & 2 & $+1$ & & 5 & 1 & $+3$ \\
\hline
\end{longtable}
\normalsize


\newpage
\begin{center}
\textbf{SPECIMEN OF EXTENDED FORM OF PELLIAN EQUATION TABLE---continued.}
\end{center}

\small
\begin{longtable}{|c|c|c|c||c|c|c|c|}
\hline
$a$ & $y$ & $x$ & $y^{2} - ax^{2}$ & $a$ & $y$ & $x$ & $y^{2} - ax^{2}$ \\
\hline
\endfirsthead
\hline
$a$ & $y$ & $x$ & $y^{2} - ax^{2}$ & $a$ & $y$ & $x$ & $y^{2} - ax^{2}$ \\
\hline
\endhead
\hline
\endfoot
23 & 4 & 1 & $+1$ & 27 & 5 & 1 & $+1$ \\
   & (3) & 3 & $-14$ &    & (5) & 1 & $-2$ \\
   & 5 & 1 & $+2$ &       & 26 & 5 & $+1$ \\
\hline
24 & 4 & 1 & $+1$ & 28 & 5 & 1 & $+1$ \\
   & (1) & 4 & $-8$ &     & (3) & 2 & $+1$ \\
   & 5 & 1 & $+1$ &       & 127 & 24 & $+1$ \\
\hline
26 & 5 & 1 & $+1$ & 29 & 5 & 1 & $+1$ \\
   & (10) & 1 & $+1$ &    & (2) & 2 & $-7$ \\
   & 51 & 10 & $+1$ &      & 70 & 13 & $-1$ \\
\hline
30 & 5 & 1 & $+1$ & 31 & 5 & 1 & $+1$ \\
   & (2) & 2 & $-5$ &     & (1) & 1 & $-6$ \\
   & 11 & 2 & $+1$ &       & 6 & 1 & $+5$ \\
   & & & &                & 1520 & 273 & $+1$ \\
\hline
32 & 5 & 1 & $+1$ & 33 & 5 & 1 & $+1$ \\
   & (1) & 1 & $-7$ &     & (2) & 1 & $-8$ \\
   & 6 & 1 & $+1$ &       & 23 & 4 & $+1$ \\
\hline
34 & 5 & 1 & $+1$ & 35 & 5 & 1 & $+1$ \\
   & (1) & 2 & $-11$ &    & (1) & 1 & $-10$ \\
   & 35 & 6 & $+1$ &       & 6 & 1 & $+1$ \\
\hline
37 & 6 & 1 & $+1$ & 38 & 6 & 1 & $+1$ \\
   & (12) & 1 & $-2$ &    & (6) & 1 & $-2$ \\
   & 73 & 12 & $+1$ &      & 37 & 6 & $+1$ \\
\hline
39 & 6 & 1 & $+1$ & 40 & 6 & 1 & $+1$ \\
   & (4) & 2 & $-3$ &     & (3) & 1 & $-4$ \\
   & 25 & 4 & $+1$ &       & 19 & 3 & $+1$ \\
\hline
41 & 6 & 1 & $+1$ & 42 & 6 & 1 & $+1$ \\
   & (2) & 2 & $+2$ &     & (2) & 1 & $-6$ \\
   & 32 & 5 & $-1$ &       & 13 & 2 & $+1$ \\
   & 2049 & 320 & $+1$ &   & & & \\
\hline
43 & 6 & 1 & $+1$ & 44 & 6 & 1 & $+1$ \\
   & (3) & 2 & $+1$ &     & (1) & 1 & $-8$ \\
   & 3482 & 531 & $+1$ &   & 199 & 30 & $+1$ \\
\hline
45 & 6 & 1 & $+1$ & 46 & 6 & 1 & $+1$ \\
   & (1) & 2 & $+1$ &     & (1) & 3 & $-5$ \\
   & 7 & 1 & $+1$ &       & 5 & 1 & $-21$ \\
   & 161 & 24 & $-2$ &     & 24335 & 3588 & $+1$ \\
   & 15129 & 2253 & $+1$ & & & & \\
\hline
47 & 6 & 1 & $+1$ & 48 & 6 & 1 & $+1$ \\
   & (5) & 1 & $-2$ &     & (1) & 1 & $-12$ \\
   & 48 & 7 & $+1$ &       & 7 & 1 & $+1$ \\
\hline
50 & 7 & 1 & $+1$ & 51 & 7 & 1 & $+1$ \\
   & (14) & 1 & $-1$ &    & (7) & 1 & $-2$ \\
   & 99 & 14 & $+1$ &      & 50 & 7 & $+1$ \\
\hline
52 & 7 & 1 & $+1$ & 53 & 7 & 1 & $+1$ \\
   & (4) & 1 & $+9$ &     & (3) & 1 & $-4$ \\
   & 649 & 90 & $+1$ &     & 182 & 25 & $+1$ \\
\hline
54 & 7 & 1 & $+1$ & 55 & 7 & 1 & $+1$ \\
   & (2) & 1 & $-5$ &     & (2) & 1 & $-6$ \\
   & 485 & 66 & $+1$ &      & 89 & 12 & $+1$ \\
\hline
56 & 7 & 1 & $+1$ & 57 & 7 & 1 & $+1$ \\
   & (2) & 1 & $-7$ &     & (1) & 1 & $-8$ \\
   & 15 & 2 & $+1$ &       & 8 & 1 & $+1$ \\
\hline
58 & 7 & 1 & $+1$ & 59 & 7 & 1 & $+1$ \\
   & (1) & 1 & $-9$ &     & (2) & 1 & $-10$ \\
   & 99 & 13 & $+3$ &      & 530 & 69 & $+1$ \\
   & 2575 & 338 & $-2$ &   & & & \\
   & 19603 & 2574 & $+1$ & & & & \\
\hline
60 & 7 & 1 & $+1$ & 61 & 7 & 1 & $+1$ \\
   & (1) & 1 & $-11$ &    & (4) & 3 & $-5$ \\
   & 8 & 1 & $+1$ &       & 39 & 5 & $+4$ \\
   & 31 & 4 & $-11$ &      & 164 & 21 & $-5$ \\
   & 231 & 30 & $+1$ &     & 29718 & 3805 & $+1$ \\
\hline
62 & 7 & 1 & $+1$ & 63 & 7 & 1 & $+1$ \\
   & (6) & 1 & $-2$ &     & (1) & 1 & $-14$ \\
   & 63 & 8 & $+1$ &       & 8 & 1 & $+1$ \\
\hline
65 & 8 & 1 & $+1$ & 66 & 8 & 1 & $+1$ \\
   & (16) & 1 & $-1$ &    & (8) & 1 & $-2$ \\
   & 129 & 16 & $+1$ &     & 65 & 8 & $+1$ \\
\hline
67 & 8 & 1 & $+1$ & 68 & 8 & 1 & $+1$ \\
   & (5) & 1 & $-3$ &     & (4) & 1 & $-4$ \\
   & 41 & 5 & $+6$ &       & 33 & 4 & $+1$ \\
   & 48842 & 5967 & $+1$ & & & & \\
\hline
69 & 8 & 1 & $+1$ & 70 & 8 & 1 & $+1$ \\
   & (3) & 1 & $-5$ &     & (2) & 1 & $-6$ \\
   & 25 & 3 & $+4$ &       & 251 & 30 & $+1$ \\
   & 7775 & 936 & $+1$ &   & & & \\
\hline
71 & 8 & 1 & $+1$ & 72 & 8 & 1 & $+1$ \\
   & (2) & 1 & $-7$ &     & (2) & 1 & $-8$ \\
   & 3480 & 413 & $+1$ &   & 17 & 2 & $+1$ \\
\hline
73 & 8 & 1 & $+1$ & 74 & 8 & 1 & $+1$ \\
   & (1) & 1 & $-9$ &     & (1) & 1 & $-10$ \\
   & 1068 & 125 & $-1$ &   & 43 & 5 & $+1$ \\
   & 2281249 & 267000 & $+1$ & & & \\
\hline
75 & 8 & 1 & $+1$ & 76 & 8 & 1 & $+1$ \\
   & (1) & 1 & $-11$ &    & (1) & 1 & $-12$ \\
   & 26 & 3 & $+1$ &       & 57799 & 6630 & $+1$ \\
   & & & &                 & & & \\
\hline
77 & 8 & 1 & $+1$ & 78 & 8 & 1 & $+1$ \\
   & (3) & 1 & $+5$ &     & (1) & 1 & $-14$ \\
   & 351 & 40 & $+1$ &     & 53 & 6 & $+1$ \\
\hline
79 & 8 & 1 & $+1$ & 80 & 8 & 1 & $+1$ \\
   & (1) & 1 & $-15$ &    & (1) & 1 & $-16$ \\
   & 80 & 9 & $+1$ &       & 9 & 1 & $+1$ \\
\hline
82 & 9 & 1 & $+1$ & 83 & 9 & 1 & $+1$ \\
   & (18) & 1 & $+1$ &    & (9) & 1 & $-2$ \\
   & 163 & 18 & $+1$ &     & 82 & 9 & $+1$ \\
\hline
84 & 9 & 1 & $+1$ & 85 & 9 & 1 & $+1$ \\
   & (6) & 1 & $-3$ &     & (4) & 1 & $-4$ \\
   & 55 & 6 & $+1$ &       & 378 & 41 & $-1$ \\
   & & & &                 & 285769 & 30996 & $+1$ \\
\hline
86 & 9 & 1 & $+1$ & 87 & 9 & 1 & $+1$ \\
   & (3) & 1 & $-5$ &     & (3) & 1 & $-6$ \\
   & 10405 & 1122 & $+1$ & & 28 & 3 & $+1$ \\
\hline
88 & 9 & 1 & $+1$ & 89 & 9 & 1 & $+1$ \\
   & (2) & 1 & $-7$ &     & (2) & 1 & $-8$ \\
   & 197 & 21 & $+1$ &     & 500 & 53 & $+1$ \\
\hline
90 & 9 & 1 & $+1$ & 91 & 9 & 1 & $+1$ \\
   & (2) & 1 & $-9$ &     & (1) & 1 & $-10$ \\
   & 19 & 2 & $+1$ &       & 1574 & 165 & $+1$ \\
\hline
92 & 9 & 1 & $+1$ & 93 & 9 & 1 & $+1$ \\
   & (1) & 1 & $-11$ &    & (1) & 1 & $-12$ \\
   & 1151 & 120 & $+1$ &   & 12151 & 1260 & $+1$ \\
\hline
94 & 9 & 1 & $+1$ & 95 & 9 & 1 & $+1$ \\
   & (1) & 1 & $-13$ &    & (1) & 1 & $-14$ \\
   & 2143295 & 221064 & $+1$ & & 39 & 4 & $+1$ \\
\hline
96 & 9 & 1 & $+1$ & 97 & 9 & 1 & $+1$ \\
   & (1) & 1 & $-15$ &    & (1) & 1 & $-16$ \\
   & 49 & 5 & $+1$ &       & 5604 & 569 & $+1$ \\
\hline
98 & 9 & 1 & $+1$ & 99 & 9 & 1 & $+1$ \\
   & (1) & 1 & $-17$ &    & (1) & 1 & $-18$ \\
   & 99 & 10 & $+1$ &      & 10 & 1 & $+1$ \\
\hline
\end{longtable}
\normalsize


\newpage
The meaning hardly requires explanation; for each number $a$, we have a series of pairs of increasing numbers, $y$, $x$, satisfying a series of equations $y^{2} = ax^{2} \pm b$; thus $a = 14$:

\begin{center}
\begin{tabular}{c|c|c}
$y$ & $x$ & $y^{2} - ax^{2}$ \\
\hline
1 & 0 & $+1$ \\
31 & 9 & $-14 \cdot 1 = -5$ \\
4 & 1 & $16 - 14 \cdot 1 = +2$ \\
11 & 3 & $121 - 14 \cdot 9 = -5$ \\
15 & 4 & $225 - 14 \cdot 16 = +1$
\end{tabular}
\end{center}

The following table, calculated under the superintendence of the Committee, extends from $a = 1001$ to $a = 1500$ (square numbers omitted); it is (with slight typographical variations) nearly but not exactly in the form of Degen's Table~I., the chief difference being that for a number $a$ having a double middle term, or of the form $\alpha^{2}+1$ (such number being further distinguished by an asterisk), the $x$, $y$ entered in the table are the solutions, not of the equation $y^{2} = ax^{2} + 1$, but of the equation $y^{2} = ax^{2} - 1$. As remarked above, if we have $y^{2} = ax^{2} - 1$, then writing $y_{1} = 2y^{2} + 1$ and $x_{1} = 2axy$, we obtain $y_{1}^{2} = ax_{1}^{2} + 1$. Moreover, for each value of $a$, in the first line, the first term, which is the integer part of $\surd a$, is separated from the other by a semicolon, and the 1, which is the corresponding first term of the second line, is omitted.

\newpage
\noindent\textbf{[948}

The calculations were made by C.~E.~Bickmore, M.A., of New College, Oxford: his values for $x$ and $y$ have been revised as presently mentioned, but it has been assumed that his values for the periods and subsidiary numbers (forming the first and second lines of each division of the table) are accurate; in fact, any error therein would cause the resulting values of $x$ and $y$ to be wildly erroneous; but (except in a single instance which was accounted for) the errors in $x$ and $y$ were in every case in a single figure or two or three figures only.

The values of $x$ and $y$ were in every case examined by substitution in the equation ($y^{2} = ax^{2} + 1$, or $y^{2} = ax^{2} - 1$, as the case may be), which should be satisfied by them. These verifications were for the most part made by A.~Graham, M.A., of the Observatory, Cambridge. As already mentioned, some errors were detected, and these have been, of course, corrected. The values of $x$, $y$ given in the table thus satisfy in every case the proper equation $y^{2} = ax^{2} + 1$, or $y^{2} = ax^{2} - 1$; on the ground above referred to, it is believed that the periods and subsidiary numbers are also accurate.

It may be remarked, in regard to the verification of the equation $y^{2} = ax^{2} \pm 1$ for large values of $a$; and $y$, it is in practice easier and safer to calculate $ax^{2} \pm 1$, and then to compare the square root thereof with the given value of $y$, than to further calculate the value of $y^{2}$.

\newpage
\begin{center}
\textbf{THE TABLE 1001 TO 1500.}
\end{center}

\small
\begin{longtable}{|c|l|c|c|}
\hline
$a$ & Continued fraction for $\surd a$ & $x$ & $y$ \\
\hline
\endfirsthead
\hline
$a$ & Continued fraction for $\surd a$ & $x$ & $y$ \\
\hline
\endhead
\hline
\endfoot
1001 & 31; 1, 1, 1, 3, 3, 2, (4) & 1060905 & 33532 \\
     & 40, 23, 35, 16, 17, 25, (13) & & \\
\hline
1002 & 31; 1, 1, 1, 8, 2, 1, 1, 1, 3, (10) & 6535248 & 2068692 \\
     & 41, 22, 39, 7, 23, 31, 26, 33, 17, (6) & & \\
\hline
1003 & 31; 1, 2, (31) & 9026 & 285 \\
     & 42, 21, (2) & & \\
\hline
1004 & 31; 1, 2, 5, 2, 2, 1, 7, 4, 1, 2, 1, 11, 1, (14) & 9633024199 & 24164859730 \\
     & 43, 20, 11, 25, 19, 41, 8, 13, 40, 17, 44, 5, 55, (4) & & \\
\hline
1005 & 31; 1, 2, 2, 1, 5, 15, 1, 2, (12) & 2950149761 & 93059568 \\
     & 44, 19, 20, 39, 11, 4, 41, 21, (5) & & \\
\hline
1006 & 31; 1, 2, 1, 1, 5, 1, 3, 2, 1, 1, 1, 1, 1, 9, 1, 20, 4, 5, 1, 1, 12, 6, 1, (30) & 2174821748 & 4544614025 \\
     & 45, 18, 29, 33, 10, 43, 15, 22, 31, 27, 30, 25, 37, 6, 55, 3, 15, 11, 30, 33, 5, 9, 53, (2) & & \\
\hline
1007 & 31; 1, 2, (1) & 476 & 15 \\
     & 46, 17, (38) & & \\
\hline
1008 & 31; 1, (2) & 127 & 4 \\
     & 47, (16) & & \\
\hline
1009* & 31; 1, (3, 3) & 540 & 17 \\
     & 48, (15, 15) & & \\
\hline
1010* & 31; 1, 3, (1, 1) & 1303 & 41 \\
     & 49, (3, 31) & & \\
\hline
1011 & 31; 1, 3, 1, (9) & 8426 & 265 \\
     & 50, 13, 47, (6) & & \\
\hline
1012 & 31; 1, 4, 3, 6, 1, 3, 8, 1, (4) & 1739932226 & 101302110 \\
     & 51, 12, 19, 9, 43, 16, 7, 48, (11) & & \\
\hline
1013* & 31; 1, 4, 1, 4, 15, 1, 2, (2, 2) & 666183931 & 12352985 \\
     & 52, 11, 44, 13, 4, 43, 19, (23, 23) & & \\
\hline
1014 & 31; 1, 5, 2, 1, 1, 1, 1, (20) & 146256965 & 5696546 \\
     & 53, 10, 23, 30, 29, 25, 38, (3) & & \\
\hline
1015 & 31; 1, 6, 10, (2) & 3528715528 & 110763 \\
     & 54, 9, 6, (29) & & \\
\hline
1016 & 31; 1, (6) & 255 & 8 \\
     & 55, (8) & & \\
\hline
1017 & 31; 1, 6, 1, 1, 5, 2, 1, 1, 1, 2, 1, (12) & 972172909322 & 29009322 \\
     & 56, 7, 8, 49, 9, 13, 41, 16, 11, 31, 32, (9) & & \\
\hline
1018* & 31; 1, 9, 1, 1, 1, 6, 2, (3, 3) & 50499870283 & 2728333 \\
     & 57, 6, 39, 23, 38, 9, 26, (17, 17) & & \\
\hline
1019 & 31; 1, 11, 1, 3, 1, 1, 1, 3, 8, 1, 5, 2, (31) & 6089923321730 & 190776436539 \\
     & 58, 5, 47, 14, 35, 25, 34, 17, 7, 49, 10, 29, (2) & & \\
\hline
1020 & 31; 1, (11) & 511 & 16 \\
     & 59, (4) & & \\
\hline
1021* & 31; 1, 20, 3, 6, 1, 3, 2, 1, 1, 12, 5, 4, 15, 1, 2, 1, 4, 1, 1, 2, 1, 1, 1, (5, 5) & 3152172803725848825030 & 986500129665669 \\
     & 60, 3, 20, 9, 44, 15, 23, 27, 36, 5, 12, 15, 4, 45, 17, 41, 12, 33, 29, 20, 33, 25, 36, (11, 11) & & \\
\hline
\end{longtable}
\normalsize


\newpage
\section*{949.}
\addcontentsline{toc}{section}{949. On Halphen's Characteristic n, in the Theory of Curves in Space}

\begin{center}
{\large\bfseries ON HALPHEN'S CHARACTERISTIC n, IN THE THEORY OF CURVES IN SPACE.}
\end{center}

\begin{flushright}
[From Crelle's Journal d.~Mathematik, t.~cxi. (1893), pp.~347---352.]
\end{flushright}

If we consider a curve in space without actual singularities, of the order (or degree) $d$, then this has a number $h$ of apparent double points (adps.), viz. taking as vertex an arbitrary point in space, we have through the curve a cone of the order $d$, with $h$ nodal lines; and Halphen denotes by $n$ the order of the cone of lowest order which passes through these $h$ lines. For a given value of $d$, $h$ is at most $=\frac{1}{2}(d-1)(d-2)$, and as shown by Halphen it is at least $=[\frac{1}{4}(d-1)^{2}]$, if we denote in this manner the integer part of $\frac{1}{4}(d-1)^{2}$. For given values of $d$, $h$, it is easy to see that $n$ must lie within certain limits, viz. if $\nu$ be the smallest number such that $\frac{1}{2}\nu(\nu+3)$ equal to or greater than $h$, then $n$ is at most $=\nu$; and moreover $n$ must have a value such that $nd$ is at least $=2h$, or say we must have $nd = 2h + \theta$, where $\theta$ is $=0$ or positive. For any given value of $d$, we thus have a finite number of forms $(d, h, n)$, and we have thus \textit{prima facie} curves in space of the several forms $(d, h, n)$: but it may very well be, and in fact Halphen finds, that when $d = 9$ or upwards, then for certain values of $h$, $n$ as above, there is not any curve $(d, h, n)$; thus $d = 9$, $h = 17$ the values of $n$ are $n = 4$, $n = 5$, but there is not any curve $d = 9$, $h = 17$ for either of these values of $n$; or say the curves $(9, 17, 4)$ and $(9, 17, 5)$ are non-existent. And in the Notes and References to the papers 302, 305 in vol.~v. of my Collected Mathematical Papers, 4to. Cambridge, 1892, see p.~615, I remarked that, in certain cases for which Halphen finds a curve $(d, h, n)$, such curve does not exist except for special configurations of the $h$ nodal lines not determined by the mere definition of $n$ as the order of the cone of lowest order which passes through the $h$ nodal lines; for instance $d = 9$, $h = 16$, $n = 4$, for which Halphen gives a curve, I find that for the existence of the curve it is not enough that the 16 lines are situate upon a quartic cone, but they must be the 16 lines of intersection of two quartic cones.

\textbf{[949}

In fact, starting from an existing curve, say the complete intersection of two given surfaces of the orders $\mu$, $\nu$ respectively, we have, as is known,
\[ d = \mu\nu, \quad 2h = \mu\nu(\mu-1)(\nu-1); \]
we find also, as will appear, $n = (\mu-1)(\nu-1)$: hence also $nd = 2h$, viz. the cone of the order $n$ through the $h$ nodal lines meets the cone of the order $d$ in these lines counting each twice, and in no other lines. I remark also that $h$ is $=\frac{1}{2}n(n+3)$ if $\mu + \nu = 4$, viz. we have here two nodal lines lying in a plane; but if $\mu + \nu > 4$, then $h$ is greater than $\frac{1}{2}n(n+3)$, viz. in this case the nodal lines are not $h$ lines taken at pleasure, but they are lines subject to the condition of lying in a cone of the order $n$. But this is not all; the $h$ nodal lines lie not only in this cone of lowest order $n$, but also in cones of the orders $n+1$, $n+2$, \ldots, $n + (\mu + \nu - 2)$ respectively: I do not for the moment attempt to determine the number of independent conditions which are hereby imposed upon the $h$ lines.

The last-mentioned theorem constitutes in fact the geometrical interpretation of results contained in Jacobi's paper ``De eliminatione variabilis e duabus aequationibus algebraicis,'' Crelle, t.~xv. (1836), pp.~101---124, [Ges.~Werke, t.~iii. pp.~295---320].

Consider the two equations
\[ U = (*\!\pitchfork\! x, y, z, w)^{\mu} = 0; \quad V = (*\!\pitchfork\! x, y, z, w)^{\nu} = 0; \]
representing surfaces of the orders $\mu$, $\nu$ respectively; since the form of the equations is quite arbitrary, we may without loss of generality assume that the point $(x, y, z) = (0, 0, 0)$ is an arbitrary point in space; and this being so, we find the equation of the cone, vertex this point, which passes through the curve of intersection of the two surfaces by the mere elimination of $w$ between the two equations. As the reasoning is exactly the same for a particular case, I write for convenience $\mu = 3$, $\nu = 4$, and consider the two equations
\[
\begin{aligned}
A_{0}w^{3} + A_{1}w^{2} + A_{2}w + A_{3} &= 0, \\
B_{0}w^{4} + B_{1}w^{3} + B_{2}w^{2} + B_{3}w + B_{4} &= 0,
\end{aligned}
\]
where the suffixes show the degrees in regard to $(x, y, z)$, viz. $A_{0}$, $B_{0}$ are mere constants, $A_{1}$, $B_{1}$ are linear functions $(*\!\pitchfork\! x, y, z)^{1}$, $A_{2}$, $B_{2}$ quadric functions $(*\!\pitchfork\! x, y, z)^{2}$, and so on. Multiplying the first equation successively by $1$, $w$, $w^{2}$, $w^{3}$, and the second equation successively by $1$, $w$, $w^{2}$, we have 7 equations from which to eliminate $1$, $w$, $w^{2}$, $w^{3}$, $w^{4}$, $w^{5}$, $w^{6}$, and the result is
\[
\begin{vmatrix}
A_{0}, & A_{1}, & A_{2}, & A_{3}, \\
& A_{0}, & A_{1}, & A_{2}, & A_{3}, \\
& & A_{0}, & A_{1}, & A_{2}, & A_{3}, \\
B_{0}, & B_{1}, & B_{2}, & B_{3}, & B_{4}, \\
& B_{0}, & B_{1}, & B_{2}, & B_{3}, & B_{4}, \\
& & B_{0}, & B_{1}, & B_{2}, & B_{3}, & B_{4}
\end{vmatrix} = 0,
\]
viz. this is an equation $\Omega = 0$ of the cone of the order $d = 12$, through the curve of intersection of the two surfaces: say this equation is $\Omega = 0$.

But if we only multiply the first equation by $1$, $w$, $w^{2}$ successively and the second equation by $1$, $w$ successively, then we have 5 equations serving to determine the ratios of $w^{5}$, $w^{4}$, $w^{3}$, $w^{2}$, $w$, $1$, viz. we have these quantities proportional to the six determinants which can be formed out of the matrix
\[
\begin{matrix}
A_{0}, & A_{1}, & A_{2}, & A_{3}, \\
& A_{0}, & A_{1}, & A_{2}, & A_{3}, \\
B_{0}, & B_{1}, & B_{2}, & B_{3}, & B_{4}, \\
& B_{0}, & B_{1}, & B_{2}, & B_{3}, & B_{4}
\end{matrix}
\]
say we have
\[ w^{5} : w^{4} : w^{3} : w^{2} : w : 1 = L : M : N : P : Q : R, \]
where $L$, $M$, $N$, $P$, $Q$, $R$ represent homogeneous functions $(*\!\pitchfork\! x, y, z)^{e}$, of the degrees $11$, $10$, $9$, $8$, $7$, $6$ respectively. We may if we please write
\[ \frac{w^{5}}{L} = \frac{w^{4}}{M} = \frac{w^{3}}{N} = \frac{w^{2}}{P} = \frac{w}{Q} = \frac{1}{R} \]
or eliminating $w$, we have the series of equations which may be written
\[
\begin{vmatrix}
L, & M, & N, & P, & Q \\
M, & N, & P, & Q, & R
\end{vmatrix} = 0,
\]
viz. we thus denote that the determinants formed with any two columns of this matrix are severally $= 0$. This of course implies that each of the determinants in question is the product of $\Omega$ and a factor which is a homogeneous function of the proper degree in $(x, y, z)$, so that the several equations are all of them satisfied if only $\Omega = 0$. We have for instance $PR - Q^{2} = A\Omega$, where $A$ is a quadric function $(*\!\pitchfork\! x, y, z)^{2}$; similarly $NR - PQ = B\Omega$, where $B$ is a cubic function $(*\!\pitchfork\! x, y, z)^{3}$, and the like as regards the other determinants.

If the ratios $x:y:z$ have any given values such that we have for these $\Omega = 0$, then $w$ has a determinate value, that is, on each line of the cone $\Omega = 0$, there is a single point of the curve of intersection of the two surfaces: the only exceptions are when, for the given values of $x:y:z$, the expressions for $w$ assume an indeterminate form, viz. $w$ has then two values, and there are upon the line two points of the curve, or what is the same thing, the line is a nodal line of the cone: the conditions for a nodal line thus are $L = 0$, $M = 0$, $N = 0$, $P = 0$, $Q = 0$, $R = 0$, viz. each of these equations is that of a cone passing through the nodal lines of the cone $\Omega = 0$; the cone of lowest order is $R = 0$, a cone of the order 6 meeting the cone $\Omega = 0$ of the order 12 in 36 lines each twice, which lines are consequently the nodal lines of the cone $\Omega = 0$. The mere condition of the 36 lines lying upon a cone of the order 6 shows that the 36 lines are not arbitrary; and we have moreover, through the 36 lines, cones of the orders 7, 8, 9, 10 and 11 respectively. Obviously the foregoing reasoning is quite general, and for the surfaces of the orders $\mu$, $\nu$ we have (as stated above) the cone $\Omega$ of the order $\mu\nu$, with $h = \frac{1}{2}\mu\nu(\mu-1)(\nu-1)$ nodal lines, the intersections (each counting twice) of this cone with a cone of the order $n = (\mu-1)(\nu-1)$; and moreover the $h$ nodal lines lie also in cones of the orders $n+1$, $n+2$, \ldots, $n+\mu+\nu-2$ respectively.

\textbf{949]}
\begin{center}
{\large\bfseries THEORY OF CURVES IN SPACE.}
\end{center}

To examine the meaning of the theorem, I form the table

\begin{center}
\begin{tabular}{c|c|c|c|l}
$\mu$, $\nu$ & $d$ & $h$ & $n$ & \multicolumn{1}{c}{$\ldots$, $n+\mu+\nu-2$} \\
\hline
2, 2 & 4 & 2 & 1, & 2, 3 \\
2, 3 & 6 & 6 & 2, & 3, 4, 5 \\
2, 4 & 8 & 12 & 3, & 4, 5, 6, 7 \\
2, 5 & 10 & 20 & 4, & 5, 6, 7, 8, 9 \\
3, 3 & 9 & 12 & 4, & 5, 6, 7 \\
3, 4 & 12 & 18 & 6, & 7, 8, 9, 10, 11 \\
3, 5 & 15 & 60 & 8, & 9, 10, 11, 12, 13, 14 \\
4, 4 & 16 & 72 & 9, & 10, 11, 12, 13, 14, 15 \\
4, 5 & 20 & 120 & 12, & 13, 14, 15, 16, 17, 18, 19
\end{tabular}
\end{center}

Here $\mu$, $\nu = 2$, 2, there are 2 nodal lines, which are arbitrary, and of course lie on cones of the orders 1, 2, 3 respectively. So $\mu$, $\nu = 2$, 3; there are 6 nodal lines, which are not arbitrary, inasmuch as they lie on a cone of the order 2; but regarding them as arbitrary lines on such a cone, we can through them draw a cone of the order 3 or any higher order, and it is thus no specialisation to say that they lie upon cones of the orders 3, 4, and 5. But going a step further $\mu = 2$, $\nu = 4$: here we have 12 nodal lines which, inasmuch as they lie on a cone of the order 3, are not arbitrary: and they are not arbitrary lines upon this cone, for they lie on a cone of the order 4, and such a cone can be drawn through at most 11 arbitrary lines on a cubic cone. In fact, upon a cone of the order $\theta$, taking at pleasure $N$ lines, the condition that it may be possible through these to draw a cone of the order $\theta+1$ is $\frac{1}{2}(\theta+1)(\theta+4) = N + 3$ at least; for if this number were $= N + 2$, then through the $N$ points we have only the improper cone $(x + \beta y + \gamma z)\, U_{\theta} = 0$, if $U_{\theta} = 0$ is the cone of the order $\theta$. It thus appears that the 12 nodal lines are not arbitrary lines on a cubic cone, but that they constitute the complete intersection of a cubic cone and a quartic cone. But through such 12 lines we may draw cones of the 5th and higher orders, and it is thus no further condition that the 12 lines lie on cones of the orders 5, 6 and 7 respectively.

\textbf{[949}

So again $\mu = 3$, $\nu = 3$; we have here 18 nodal lines which, inasmuch as they lie on a cone of the order 4, are not arbitrary: and they are not arbitrary lines on this cone inasmuch as they lie also on a cone of the order 5, and such a cone can be drawn through at most 17 arbitrary lines on the quartic cone: it thus appears that the 18 nodal lines are 18 out of the 20 lines of intersection of a quartic cone and a quintic cone. But there is no further condition, for through such lines we can draw a cone of the order 6 or any higher order and thus the lines lie on cones of the orders 6, 7 and 8 respectively. It appears probable however that, for higher values of $\mu$, $\nu$, it would be necessary to take account not only (as in these examples) of the cones of the orders $n$ and $n+1$, but of those of higher orders $n+2$, \&c.; and thus that it is not the true form of the theorem to say that the $h$ nodal lines must be $h$ out of the $n(n+1)$ lines of intersection of two cones of the orders $n$ and $n+1$ respectively.

It appears, by what precedes, that the $h = \frac{1}{2}\mu\nu(\mu-1)(\nu-1)$, lines which are the nodal lines of the cone of arbitrary vertex which passes through the curve of intersection of two surfaces of the orders $\mu$, $\nu$ respectively, form a remarkable special system of lines, which well deserve further study. I remark also that, without having proved the negative, it seems to me clear that given the values of $d$, $h$, $n$ it is only in the cases where the $h$ lines form some such special system (and not in the general case where the $h$ nodal lines are any lines whatever on a cone of the order $n$) that there exists a curve $(d, h, n)$; and thus that the question for further investigation is, for given values of $(d, h, n)$ to determine the conditions necessary for the existence of a curve in space with these characteristics $(d, h, n)$.


\newpage
\section*{950.}
\addcontentsline{toc}{section}{950. On the Sextic Resolvent Equations of Jacobi and Kronecker}

\begin{center}
{\large\bfseries ON THE SEXTIC RESOLVENT EQUATIONS OF JACOBI AND KRONECKER.}
\end{center}

\begin{flushright}
[From Crelle's Journal d.~Mathematik, t.~cxm. (1894), pp.~42---49.]
\end{flushright}

The equations referred to are: the first of them, that given by Jacobi in the paper ``Observatiunculae ad theoriam aequationum pertinentes,'' Crelle, t.~xiii. (1835), pp.~340---352, [Ges.~Werke, t.~iii., pp.~269---284], under the heading ``Observatio de aequatione sexti gradus ad quam aequationes sexti gradus revocari possunt,'' and the second, that of Kronecker in the note ``Sur la resolution de l'\'{e}quation du cinqui\`{e}me degr\'{e},'' Comptes Rendus, t.~XLVI. (1858), pp.~1150---1152. Jacobi's equation is closely connected with that obtained by Malfatti in 1771, see Brioschi's paper ``Sulla resolvente di Malfatti per l'equazione del quinto grado,'' Mem.~R.~Ist.~Lomb., t.~ix. (1863); but the characteristic property first presents itself in Jacobi's form, and I think the equation is properly described as Jacobi's resolvent equation. The other equation has been always known as Kronecker's resolvent equation; it belongs to the class of equations for the multiplier of an elliptic function considered by Jacobi in the paper ``Suite des notices sur les fonctions elliptiques,'' Crelle, t.~iii. (1828), pp.~303---310, see p.~308, [Ges.~Werke, t.~i., pp.~255---263, see p.~261]: say Kronecker's equation belongs to the class of Jacobi's Multiplier Equations. We have in regard to it the paper by Brioschi, ``Sul metodo di Kronecker per la resoluzione delle equazioni di quinto grado,'' Atti Ist.~Lomb., t.~i. (1858), pp.~275---282, and see also the ``Appendice terza'' to his translation of my Elliptic Functions (Milan, 1880): it seems to me however that the theory of Kronecker's equation has not hitherto been exhibited in the clearest form.

I consider the forms

\begin{center}
\begin{tabular}{c@{\hspace{2em}}c}
12345 & 13524 \\
13254 & 12435 \\
24315 & 23541 \\
35421 & 34152 \\
41532 & 45213 \\
52143 & 51324
\end{tabular}
\end{center}

which, if in the first instance the figures are regarded as points, represent the twelve pentagons which can be formed with the points 1, 2, 3, 4, 5; each form in the right-hand column is derived from the corresponding form in the left-hand column by stellation, say we have $13524 = S\,12345$, and so in other cases.

A pentagon is in general reversible, but we sometimes consider it as irreversible (viz. we distinguish between the pentagons 12345 and 15432); when this is so, we write $15432 = R\,12345$, and we have thus twelve new forms, in all twenty-four forms. The symbols $R$, $S$ are such that $R^{2} = 1$, $S^{2} = R$, $RS = SR$, $S^{4} = 1$. But for a reversible pentagon, there is no occasion to use the symbol $R$, and we have simply $S^{2} = 1$.

In a somewhat different point of view, we may for an irreversible pentagon write

$12345$ to denote any one of the forms $12345$, $23451$, $34512$, $45123$, $51234$,

$R\,12345$ \quad\; to denote any one of the forms $15432$, $54321$, $43215$, $32154$, $21543$,

and for a reversible pentagon $12345$ to denote any one of these same ten forms.

Each pentagon gives thus ten forms, viz. we have in all 120 forms which all are the different arrangements of the five figures. But further regarding the arrangement 12345 as positive, then the forms in the left-hand column are each of them positive, and the ten forms derived from any one of these are each of them positive, that is, the forms in the left-hand column give all the 60 positive arrangements of the five figures: and similarly the forms in the right-hand column give all the 60 negative arrangements of the five figures.

Taking 1, 2, 3, 4, 5 to denote any quantities, or say the five roots $x_{1}$, $x_{2}$, $x_{3}$, $x_{4}$, $x_{5}$ of a quintic equation, we regard $12345$, \ldots, as denoting functions of these roots: in particular, $12345$ may denote a cyclic reversible function, the analogue of the reversible pentagon, or it may denote a cyclic irreversible function, the analogue of the irreversible pentagon.

\textbf{Jacobi's resolvent equation.}

The most simple course is to take $12345$ a cyclic reversible function of the roots $x$; a root of the resolvent equation is then $12345 - 13524$, $= (1-S)\,12345$, and the six roots are
\[
\begin{aligned}
z_{1} &= (1-S)\,12345, \\
z_{2} &= (1-S)\,13254, \\
z_{3} &= (1-S)\,24315, \\
z_{4} &= (1-S)\,35421, \\
z_{5} &= (1-S)\,41532, \\
z_{6} &= (1-S)\,52143.
\end{aligned}
\]
Here effecting on the roots $x$ any positive substitution whatever, we permute inter se the roots $z$; but effecting on the roots $x$ any negative substitution whatever, then the roots $z$ are permuted inter se, but at the same time each root $z$ changes its sign: hence the coefficients of the equation having the roots $z$ are not absolutely unaltered, but they are each of them multiplied by a determinate factor (1 or $-1$) according as the substitution is positive or negative; such an equation is not properly invariantive in regard to the general substitutions of the roots, but only in regard to the positive substitutions: or say it is a \textit{semi-invariant} equation. But observe that the equation having the roots $z^{2}$ is invariantive as well in regard to the positive as to the negative substitutions; hence if from the equation in $z$ we eliminate $z$ by means of the equation $y = z^{2}$, we obtain an invariant equation of the order 6 in $y$.

\vspace{1em}
\hrule
\vspace{1em}

\begin{center}
{\small\textit{End of Pages 451--500}}
\end{center}

\end{document}
